\section{Defining a quiver in Lean}\label{sec:11-path-algebras:03-defining-a-quiver:1}

A quiver is encoded as a \texttt{structure}, parameterized by a type of vertices and a type of arrows:

\begin{lstlisting}
structure Quiver (V : Type) (A : Type) where
  source : A → V
  target : A → V
\end{lstlisting}

This is a direct transcription: \texttt{V} is the type of vertices, \texttt{A} the type of arrows, and \texttt{source}/\texttt{target} are exactly \(s\) and \(t\) above.

\begin{mathreading}

\texttt{Quiver V A} is a quiver \(Q = (V, A, s, t)\): the data of two functions \(s, t : A \to V\), i.e.~a parallel pair of arrows \(A \rightrightarrows V\) in \(\mathbf{Set}\). A quiver is precisely the underlying data of a small category \emph{before} composition and identities are imposed: the raw generators from which the free category (paths) is built.

\end{mathreading}

\subsection{The example quiver, formalized}\label{sec:11-path-algebras:03-defining-a-quiver:2}

The three vertices are encoded as \texttt{Fin 3} (values \texttt{0, 1, 2}, standing for vertices \texttt{1, 2, 3}) and the two nontrivial arrows as an inductive type:

\begin{lstlisting}
inductive ExampleArrow where
  | alpha : ExampleArrow   -- 0 → 1
  | beta : ExampleArrow     -- 1 → 2

def exampleQuiver : Quiver (Fin 3) ExampleArrow where
  source := fun a => match a with
    | ExampleArrow.alpha => 0
    | ExampleArrow.beta => 1
  target := fun a => match a with
    | ExampleArrow.alpha => 1
    | ExampleArrow.beta => 2
\end{lstlisting}

\texttt{inductive ExampleArrow where | alpha : ExampleArrow | beta : ExampleArrow} defines a brand-new type with exactly two values, \texttt{alpha} and \texttt{beta}. This is the same \texttt{inductive} mechanism that builds \texttt{Nat} (Chapter 1) and \texttt{Bool}, just with two constructors that carry no extra data.

\begin{mathreading}

This is the concrete quiver \(Q\) with vertex set \(V = \{0,1,2\}\) (encoded as \(\mathrm{Fin}\,3\)) and arrow set \(A =
\{\alpha, \beta\}\), where \(s(\alpha)=0,\ t(\alpha)=1\) and \(s(\beta)=1,\
t(\beta)=2\): the linear \(A_3\) quiver, drawn exactly as on paper:

\end{mathreading}

\begin{center}
% The running example quiver: 0 -> 1 -> 2, arrows alpha and beta.
% Source: 11-path-algebras/03-defining-a-quiver.md.
\begin{tikzcd}
  0 \arrow[r, "\alpha"] & 1 \arrow[r, "\beta"] & 2
\end{tikzcd}

\end{center}

\begin{tabular}{ll}
\toprule
Symbol & Lean \\
\midrule
\bottomrule
\(0, 1, 2 \in V\) (``vertices'') & \texttt{(0 : Fin 3)}, \texttt{(1 : Fin 3)}, \texttt{(2 : Fin 3)} \\
\(\alpha, \beta \in A\) (``arrows'') & \texttt{ExampleArrow.alpha}, \texttt{ExampleArrow.beta} \\
\(s, t : A \to V\) (``source, target'') & \texttt{exampleQuiver.source}, \texttt{exampleQuiver.target} \\
\end{tabular}


The two-element inductive type is the finite set \(A = \{\alpha,\beta\}\), and the \texttt{match}-defined \texttt{source}/\texttt{target} are the functions \(s, t : A \to
V\) given by their value tables: \texttt{source alpha = 0}, \texttt{target alpha = 1}, \texttt{source beta = 1}, \texttt{target beta = 2}. This is exactly the arrowheads and tails drawn above.

\textbf{Mathlib equivalent.} Mathlib's own \href{https://loogle.lean-lang.org/?q=Quiver}{\texttt{Quiver}} class (the one this chapter's opening section mentioned building from scratch instead of reusing) encodes arrows differently. Instead of one flat arrow type \texttt{A} plus separate \texttt{source}/\texttt{target : A → V} functions, it bakes the endpoints into the arrow's \emph{type} directly, \texttt{Hom : V → V → Sort*} (with notation \texttt{a ⟶ b}). Thus an arrow from \texttt{i} to \texttt{j} simply \emph{has type} \texttt{i ⟶ j}:

\begin{lstlisting}
inductive MyArrow : Fin 3 → Fin 3 → Type
  | alpha : MyArrow 0 1
  | beta : MyArrow 1 2

instance : Quiver (Fin 3) := ⟨MyArrow⟩
\end{lstlisting}

\texttt{MyArrow 0 1} has exactly the one constructor \texttt{alpha}. \texttt{MyArrow i j} for any other pair \texttt{(i, j)}, in particular the ``backwards'' or ``no such arrow'' cases, has no constructors at all, and is thus simply an empty type. There is no \texttt{source}/\texttt{target} to state or prove separately, and no \texttt{h : Q.source a = v} side-condition to discharge with \texttt{rfl} later (Chapter 11 §4). An ill-typed composition is rejected by the type checker before a proof obligation is even reached, one step earlier than the book's own encoding catches the same mistake.
